\documentclass[12pt,a4paper]{article}

\usepackage[margin=1in]{geometry}
\usepackage{setspace}
\usepackage{titlesec}
\usepackage{graphicx}
\usepackage{hyperref}
\usepackage{enumitem}

\setstretch{1.3}

\titleformat{\section}{\large\bfseries}{\thesection}{1em}{}
\titleformat{\subsection}{\normalsize\bfseries}{\thesubsection}{1em}{}

\title{\textbf{CSE400 - Fundamentals of Probability in computing}\\Lecture 3 \\Lecture scribe }
\author{}
\date{}

\begin{document}

\maketitle
\thispagestyle{empty}

\tableofcontents
\newpage

\section{Introduction}

The third lecture of the CSE400 course focuses on the initiation and execution of the course project. The project component is an important part of the course and is designed to help students apply theoretical concepts to real-world problems. Through this project, students are expected to engage in problem identification, mathematical modeling, coding, and analytical evaluation.

The lecture primarily introduces the structure of the project, its milestones, the types of submissions required, and the evaluation process. It also explains the artifacts students must create during the project lifecycle, including concept evolution maps, scribes, and multimedia artifacts.

The goal of the project is to ensure that students not only understand theoretical concepts but also develop the ability to apply them in a structured and research-oriented manner.

\section{Understanding the CSE400 Course Project}

The CSE400 course includes a project component that encourages collaborative work, analytical thinking, and practical implementation of concepts learned during the course.

\subsection{Project Weightage}

The project contributes \textbf{30\% of the total course evaluation}. This makes the project a significant part of the course assessment.

\subsection{Team Formation}

Students are required to form teams to work on the project. Team collaboration plays a crucial role in completing the various milestones and deliverables associated with the project.

\subsection{Deadline for Team Formation}

The deadline for forming teams is:

\begin{center}
\textbf{17 January 2026 (Saturday – End of Day)}
\end{center}

After teams are formed, students begin working on the project according to the milestone structure defined in the course.

\section{Project Execution Guidelines}

The project is carried out through a series of milestones and structured submissions.

\subsection{Milestones}

The project consists of \textbf{six milestones}, labeled from \textbf{M1 to M6}. Each milestone represents a specific phase in the project development process.

Each team is required to submit \textbf{one submission per milestone}.

\subsection{Artifacts to be Produced}

During the project, students must produce several types of artifacts, including:

\begin{itemize}
\item Concept Evolution Maps
\item Scribe Documentation
\item Multimodal Artifacts
\item Question-Driven Artifacts
\item Collaboration and Team Dynamics Artifacts
\end{itemize}

These artifacts document the progress and evolution of the project throughout the semester.

\subsection{Deliverables}

The project deliverables may include:

\begin{itemize}
\item Code implementations
\item Reports
\item Videos
\item Other artifacts specified during the course
\end{itemize}

\subsection{Project Assessment}

The evaluation of the project includes several components:

\begin{itemize}
\item Team assessment before the mid-semester examination
\item Team assessment after the mid-semester examination
\item Final project viva
\item Final submission at the end of the course
\end{itemize}

These assessments help measure the progress and understanding of the project work.

\section{Project Milestones (M1 – M6)}

The project development process is divided into six milestones, each focusing on a particular stage of the project.

\subsection{Milestone 1 (M1): Kickstart}

The first milestone focuses on initiating the project and defining the problem.

\textbf{Activities involved in this milestone include:}

\begin{enumerate}
\item Team Formation
\item Area Identification
\item Background Study
\item Motivation
\item Problem Formulation
\end{enumerate}

Students first form teams and then identify a suitable area for their project. After selecting the area, they conduct a background study to understand existing knowledge related to the problem. The motivation behind choosing the problem must also be clearly explained.

Finally, the problem is formulated in a clear and structured way so that it can be analyzed and solved in later stages.

\subsection{Milestone 2 (M2): Mathematical Modeling}

In this milestone, the real-world problem identified earlier is converted into a mathematical representation.

The mathematical model may involve several probabilistic concepts, including:

\begin{itemize}
\item Random Variables (RV)
\item Probability Mass Function (PMF)
\item Probability Density Function (PDF)
\item Cumulative Distribution Function (CDF)
\item Multivariate Random Variables
\item Joint PMF
\item Joint PDF
\item Joint CDF
\end{itemize}

Mathematical modeling helps represent complex real-world problems in a structured mathematical form. This makes it possible to analyze the problem using computational techniques.

\subsection{Milestone 3 (M3): Coding}

After mathematical modeling, the next stage involves implementing the model using programming.

The tasks in this milestone include:

\begin{itemize}
\item Implementing the mathematical model using code
\item Performing simulations
\item Conducting computational analysis
\end{itemize}

The purpose of this stage is to verify whether the mathematical model behaves as expected when implemented computationally.

\subsection{Milestone 4 (M4): Inference}

This milestone focuses on the use of randomized algorithms.

Students are required to:

\begin{enumerate}
\item Select a randomized algorithm
\item Understand how the algorithm works
\item Implement the algorithm using code
\end{enumerate}

Randomized algorithms incorporate randomness as part of their logic. These algorithms can sometimes provide efficient solutions to problems that are difficult to solve deterministically.

\subsection{Milestone 5 (M5): Randomized Algorithms}

After implementing the randomized algorithm, students apply it to their project problem.

The main tasks include:

\begin{itemize}
\item Applying the randomized algorithm to the problem domain
\item Analyzing the results produced by the algorithm
\end{itemize}

Students may also compare the performance of randomized algorithms with deterministic algorithms to observe differences and gain insights.

\subsection{Milestone 6 (M6): Bounds and Analysis}

The final milestone focuses on theoretical analysis and final submission.

Activities in this milestone include:

\begin{itemize}
\item Deriving theoretical bounds
\item Analyzing the results obtained from the algorithms
\item Compiling the findings
\item Preparing the final submission
\end{itemize}

The final submission may include:

\begin{itemize}
\item Code
\item Reports
\item Videos
\item Supporting artifacts
\end{itemize}

\section{Submission 1 – Concept Evolution Maps}

Concept evolution maps are used to visualize how the project idea develops over time.

Students may use tools such as:

\begin{itemize}
\item Miro Concept Map Tool
\item Draw.io (diagrams.net)
\end{itemize}

These maps help represent the relationship between different ideas and show how the project concept evolves throughout the semester.

\section{Submission 2 – Scribe (Learning Reflection Logs)}

Students are required to prepare scribes documenting lecture learnings and project-related decisions.

There are two types of scribes.

\subsection{Lecture Scribe}

Lecture scribes summarize the content discussed during lectures.

\begin{itemize}
\item Two groups are assigned to prepare a lecture scribe.
\item The scribe should clearly explain the lecture content.
\item The minimum length required is \textbf{8–10 pages}.
\end{itemize}

\subsection{Project Scribe}

Project scribes document the decision-making process during project development.

They include:

\begin{itemize}
\item Decision logs
\item Explanation of why one option was chosen over another
\item Constraints faced
\item Alternatives considered
\item Final decisions taken
\item Supporting evidence
\end{itemize}

\subsection{Trade-Off Analysis}

Students may use trade-off matrices to compare different alternatives based on factors such as:

\begin{itemize}
\item Cost
\item Performance
\item Risk
\end{itemize}

These comparisons help justify the decisions made during the project.

\section{Submission 3 – Multimodal Artifacts}

Students must also produce multimedia artifacts to explain their project work.

\subsection{Video Requirements}

\begin{itemize}
\item One video must be submitted for each milestone.
\item The video duration should be \textbf{10–15 minutes}.
\end{itemize}

\subsection{Video Content}

The video should explain:

\begin{itemize}
\item The work completed in the milestone
\item Coding and simulation results if applicable
\end{itemize}

Students may use tools such as PowerPoint, Google Slides, or screen recording software to create these videos.

\subsection{Types of Multimedia Artifacts}

Examples include:

\begin{itemize}
\item Think-Aloud Videos
\item One-Minute Insight Videos
\item Project Demonstration Videos
\end{itemize}

The evaluation focuses mainly on the quality of the content rather than video editing skills.

\section{Summary}

Lecture 3 introduces the overall framework for executing the CSE400 course project. The lecture explains the project structure, milestone-based execution process, required artifacts, and evaluation criteria.

The project contributes 30\% to the overall course grade and is divided into six milestones that guide students through problem formulation, mathematical modeling, coding, algorithm implementation, and analysis.

Students are also required to prepare concept evolution maps, lecture and project scribes, and multimedia artifacts to document their learning and project development process.

Overall, the lecture emphasizes a structured approach to project execution, ensuring that students apply theoretical concepts in practical and analytical contexts.

\end{document}